\documentclass[uplatex,a4j,12pt,fleqn,dvipdfmx]{ujarticle}

% 2. 基本グラフィック・色・レイアウト
\usepackage{graphicx,color}
\usepackage[margin=20mm]{geometry}
\usepackage{multicol,varwidth,float,caption}

% 3. emathコアパッケージ
\usepackage[deepenum,notMy]{emath}

% emath拡張パッケージ群
\usepackage{hako}        % 穴埋め（箱）
\usepackage{EMfbox}      % 囲み枠
\usepackage{abcgreek}    % ギリシャ文字
\usepackage{emathEy}     % 記号・特殊文字
\usepackage{emathG}      % 幾何（図形）
\usepackage{emathT}      % 表作成
\usepackage{emathPs}     % PostScript連携
\usepackage{emathR}      % 実数計算
\usepackage{emathAe}     % 解答作成
\usepackage{emathThm}    % 定理環境
\usepackage{EMpict2e}    % 描画拡張
\usepackage[deluxe]{emathOtf} % OTFフォント（太字などの対応）

% 5. 各種定義と書式設定
% 長丸囲み文字の定義
\def\nagamaru#1{\mbox{\pIIebNagamaru{#1}}}

% 定理環境（「仮定」など）の共通テスト風デザイン
\theorembodyfont{\normalfont}
\theoremstyle{boxed}
\newtheorem<frame=rectbox,frameoption={oval=8pt,preitem=~,postitem=~}>{katei}{仮定}

% キャプションの日本語化（図1: → 図 1）
\captionsetup{figurename=図,labelsep=quad}

% カウンターのリセット設定（大問ごとに数式や図の番号をリセット）
\resetcounter{equation}[enumi]
\resetcounter{figure}[enumi]
\resetcounter{table}[enumi]
\EMfigtrue
\apnlist{\setlength{\listparindent}{1zw}}

% 大問ラベルの書式定義（専用の箇条書き環境として定義）
\newenvironment{QuestionEnum}{%
    \begin{enumerate}<apnlist={\leftmargin1zw\itemindent4zw}>'{\bfseries\Large 第\arabic{enumi}問}'%
}{%
    \end{enumerate}%
}
\begin{document}
\abovedisplayskip=7.0pt plus 2.0pt minus 5.0pt
\belowdisplayskip=7.0pt plus 2.0pt minus 5.0pt

\openKaiFile

% ========== 冊子全体表紙 ==========
\begin{center}
    {\bfseries\Huge Arcaea共通テスト} \bigskip\par
    \Large （全　問　必　答） \bigskip\par
    \centerline{2027/01/16}
\end{center}
\newpage

\section*{問題}

% ========== 科目：Arcaea ==========
\subsection*{Arcaea}
\ReadTeXFile{./arcaea/2027100}
\newpage

\begin{QuestionEnum} % ← ここを変更
    \item \kaitou\ReadTeXFile{./arcaea/2027101}
    \newpage
    \item \kaitou\ReadTeXFile{./arcaea/2027102}
    \newpage
    \item \kaitou\ReadTeXFile{./arcaea/2027103}
    \newpage
\end{QuestionEnum} % ← ここを変更

% ========== 科目：文科 ==========
\subsection*{文科}
\ReadTeXFile{./arts/2027200}
\newpage

\begin{QuestionEnum} % ← ここを変更
    \item \kaitou\ReadTeXFile{./arts/2027201}
    \newpage
    \item \kaitou\ReadTeXFile{./arts/2027202}
    \newpage
    \item \kaitou\ReadTeXFile{./arts/2027203}
    \newpage
\end{QuestionEnum} % ← ここを変更

% ========== 科目：理科 ==========
\subsection*{理科}
\ReadTeXFile{./science/2027300} 
\newpage

\begin{QuestionEnum} % ← ここを変更
    \item \kaitou\ReadTeXFile{./science/2027301}
    \newpage
    \item \kaitou\ReadTeXFile{./science/2027302}
    \newpage
    \item \kaitou\ReadTeXFile{./science/2027303}
    \newpage
\end{QuestionEnum} % ← ここを変更

\closeKaiFile

% ========== 解答・解説セクション ==========
% \newpage
% \section*{正答・配点}
% \enlargethispage{2\baselineskip}
% \includegraphics*[scale=.45]{fig/KT1Aans.pdf}

\newpage
% \section*{略解}
% \inputKaiFile

\end{document}